 \documentclass[11pt]{article}
\usepackage{amsmath,amsthm}
\usepackage{named}
\usepackage{changepage}

\newtheorem{proposition}{Proposition}
\newtheorem{lemma}{Lemma}

\def\o{\overline}
\def\bs{\backslash}
\def\std{\mathit{Std}}
\def\head{\mathit{Head}}
\def\body{\mathit{Body}}
\def\p2f{\mathit{p2f}}
\def\ph{\mathit{PH}}
\def\asm{\mathit{Asm}}
\def\ug{\mathit{UG}}
\def\HT{\mathit{HT}}
\def\fnt{\mathit{finite}}
\def\gringo{{\sc gringo}}
\def\clingo{{\sc clingo}}
\def\mg{mini-\gringo}
\def\ar{\leftarrow}
\def\rar{\rightarrow}
\def\uc{\widehat\forall}
\def\lrar{\leftrightarrow}
\def\beq{\begin{equation}}
\def\eeq#1{\label{#1}\end{equation}}
\def\ba{\begin{array}}
\def\ea{\end{array}}
\def\kbugsd{\emph{KB, UG, SD}}
\newcommand{\co}{\ensuremath{\mathit{count}}}
\newcommand{\no}{\ensuremath{\mathit{not}}}

\title{Reasoning about Logic Programming Solutions\\
  to the N-Queens Puzzle\\
  (Draft)}
\author{Vladimir Lifschitz\\ University of Texas}

\begin{document}
% \date{}
\maketitle

\begin{adjustwidth}{-1cm}{-1cm}

\section{Introduction}
  
We would like to show how reasoning in classical first-order logic can
be used to verify equivalence of two encodings of the $n$-Queens
puzzle \cite[Tables~2 and~3]{geb11b}:
\medskip

\begin{center}Program $A$\end{center}
$$
\ba l
n\,\{ q(1\,..\,n,1\,..\,n) \}\, n\\
\ar q(X,Y_1) \land q(X,Y_2) \land Y_1 < Y_2\\
\ar q(X_1,Y) \land q(X_2,Y) \land X_1 < X_2\\
\ar q(X_1,Y_1) \land q(X_2,Y_2) \land X_1 < X_2 \land |X_1-X_2| = |Y_1-Y_2|
\ea
$$

\medskip
\begin{center}Program $B$\end{center}
$$\ba l
\{ q(1\,..\,n,1\,..\,n) \}\\
\ar X = 1\,..\,n \land \no\ 1\, \{ q(X,Y) \}\, 1\\
\ar Y = 1\,..\,n \land \no\ 1\, \{ q(X,Y) \}\, 1\\
\ar Z = 1\,..\,2\times n-1 \land \no\   \{ q(X,Y) : Z = X-Y+n \}\, 1\\
\ar Z = 1\,..\,2\times n-1 \land \no\   \{ q(X,Y) : Z = X+Y-1 \}\, 1\\
\ea
$$

\medskip
Some rules in these programs implicitly refer to the \co\ aggregate, and
we view choice rules as abbreviations for rules with an atom in
the head.  Written out in full, the programs become
\begin{align}
& q(X,Y) \ar \no\   \no\   q(X,Y) \land X = 1\,..\,n \land Y = 1\,..\,n	\label{a1}\\
& \ar \no\ n\leq  \co\{ X,Y : q(X,Y) \land X = 1\,..\,n \land Y = 1\,..\,n \}
  \leq n
  \label{a2}\\
& \ar q(X,Y_1) \land q(X,Y_2) \land Y_1 < Y_2\label{a3}\\
& \ar q(X_1,Y) \land q(X_2,Y) \land X_1 < X_2\label{a4}\\
& \ar q(X_1,Y_1) \land q(X_2,Y_2) \land X_1 < X_2 \land |X_1-X_2| = |Y_1-Y_2|
  \label{a5}
\end{align}

\noindent and

\begin{align}
& q(X,Y) \ar \no\   \no\   q(X,Y) \land X = 1\,..\,n \land Y = 1\,..\,n\label{b1}\\
& \ar X = 1\,..\,n \land \no\ 1\leq \co\{ Y : q(X,Y) \} \leq 1\label{b2}\\
& \ar Y = 1\,..\,n \land \no\ 1\leq \co\{ X : q(X,Y) \} \leq 1\label{b3}\\
& \ar Z = 1\,..\,2\times n-1
  \land \no\   \co\{ X,Y : q(X,Y) \land Z = X-Y+n \} \leq  1
  \label{b4}\\
& \ar Z = 1\,..\,2\times n-1
  \land \no\   \co\{ X,Y : q(X,Y) \land Z = X+Y-1 \} \leq  1
  \label{b5}
\end{align}

\section{The target Language}\label{sec:target}

For any dialect~$D$ of \mg, $\sigma^\#(D)$ stands for the signature
obtained from the signature~$\sigma(D)$ \cite[Section~3.1]{fan26} by adding
\begin{itemize}
\item the sort \emph{set},
\item the binary predicate symbol~$\in$, with the first argument of the sort
  \emph{general} and the second argument of the sort \emph{set}, and
\item the unary function symbol~\co, with the argument of the
  sort~\emph{set} and the value of the sort \emph{general}.
\end{itemize}
The symbols $\not\in$ and $\subseteq$ are understood as abbreviations.

The use of set variables for representing aggregates by first-order formulas
was proposed by \citeauthor{fan22a} [\citeyear{fan22a}].  In this paper, set
variables are meant to range over arbitrary sets of precomputed terms.
The expression $\co(S)$ is assumed to denote the numeral corresponding to the
cardinality of~$S$ if~$S$ is finite, and some non-numeral otherwise.  Thus
the formula $\exists I(I=\co(S))$, where~$I$ is an integer variable, expresses
that~$S$ is finite.  This formula will be denoted by $\fnt(S)$.  When we
reason about a user guide \cite[Section~3]{fan23a}, we may wish to assume
that the extents of the input predicates are finite.

We will need to extend the signature $\sigma^\#(D)$ by object constants of
the sort \emph{set} and by function constants with values of the sort
\emph{set}.  The intended use of these additional constants will be
described by explicit definitions, such as
$$\forall V(V\in s \lrar F(V))$$
and
$$\forall IV(V\in s(I) \lrar F(I,V)).$$

\section{Representing rules by formulas}
\label{sec:repr}

To represent rules of Programs $A$ and $B$ by first-order formulas, we take~$D$
to be the dialect obtained from the
dialect~$D_0$ \cite[Section~2]{fan26} by adding the absolute value symbol
(it is used in rule~(\ref{a5})).

To transform rule~(\ref{a1}) into a formula, we rewrite it in numeric normal
form
$$q(X,Y) \ar \no\   \no\   q(X,Y) \land X = 1\,..\,U \land Y = 1\,..\,U
\land U=n$$
\cite[Section~4]{fan26} and then apply the
transformation~$\nu$ \cite[Section~5]{fan26}:
\beq
\forall I J N(\neg\neg q(I,J) \land 1\leq I \leq N \land 1 \leq J \leq N
\land N=n \to q(I,J)).
\eeq{a1f}

The representation of rule~(\ref{a2}) includes the constant~$s_1$, which
is meant to be used in accordance with the definition
\beq
\forall V(V\in s_1 \lrar \exists IJ(V=(I,J)\land q(I,J) \land 1\leq I\leq n
\land 1\leq J\leq n)).
\eeq{d1}
The formula corresponding to rule~(\ref{a2}) is
\beq
\neg\neg(n\leq \co(s_1)\leq n).
\eeq{a2f}

Rules~(\ref{a3})--(\ref{a5}) become
\beq
\forall XY_1Y_2\neg(q(X,Y_1) \land q(X,Y_2) \land Y_1 < Y_2),
\eeq{a3f}
\beq
\forall X_1X_2Y\neg(q(X_1,Y) \land q(X_2,Y) \land X_1 < X_2)
\eeq{a4f}
and
\beq
\forall I_1I_2J_1J_2
\neg(q(I_1,J_1) \land q(I_2,J_2) \land I_1<I_2 \land |I_1-I_2|=|J_1-J_2|).
\eeq{a5f}

Rule~(\ref{b1}) is the same as~(\ref{a1}).  The representation
of rule~(\ref{b2}) includes the function constant~$s_2$, which
is meant to be used in accordance with the definition
\beq
\forall IV(V\in s_2(I) \lrar q(I,V)).
\eeq{d2}
The formula corresponding to rule~(\ref{b2}) is
\beq
\forall I\neg(1\leq I\leq n \land \neg (1\leq \co(s_2(I))\leq 1)).
\eeq{b2f}
Translating rule~(\ref{b3}) is similar; it involves the formulas
\beq
\forall IV(V\in s_3(I) \lrar q(V,I))
\eeq{d3}
and 
\beq
\forall I\neg(1\leq I\leq n \land \neg (1\leq \co(s_3(I))\leq 1)).
\eeq{b3f}

To transform rules~(\ref{b4}) and~(\ref{b5}) into formulas, we rewrite them
first in numeric normal form:
% \begin{adjustwidth}{-8mm}{-8mm}
$$%  \beq
\ba l
\ar Z = 1\,..\,2\times U-1
\land \no\   \co\{ X,Y : q(X,Y) \land Z = X-Y+U \} \leq 1 \land U=n,
\\
% \eeq{b4nnf}
%%\beq
\ar Z = 1\,..\,2\times U-1
  \land \no\   \co\{ X,Y : q(X,Y) \land Z = X+Y-1 \} \leq  1 \land U=n.
  \ea
  $$%\eeq{b5nnf}
% \end{adjustwidth}
These rules are converted into formulas containing new constants~$s_4$
and~$s_5$:
\beq
\forall KNV(V\in s_4(K,N) \lrar \exists IJ(V=(I,J) \land q(I,J)
   \land K=I-J+N)),
\eeq{d4}
\beq
\forall KN\neg(1\leq K\leq 2\times N-1 \land\neg\co(s_4(K,N))\leq 1 \land N=n),
\eeq{b4f}
\beq
\forall KV(V\in s_5(K) \lrar \exists IJ(V=(I,J) \land q(I,J)
   \land K=I+J-1)),
\eeq{d5}
\beq
\forall KN\neg(1\leq K\leq 2\times N-1 \land\neg\co(s_5(K))\leq 1 \land N=n).
\eeq{b5f}

To sum up,
\begin{itemize}
\item rules of Program~$A$ are represented by formulas~(\ref{a1f}),
  (\ref{a2f}), (\ref{a3f}), (\ref{a4f}), (\ref{a5f});
\item rules of Program~$B$ are represented by formulas~(\ref{a1f}),
 (\ref{b2f}), (\ref{b3f}), (\ref{b4f}), (\ref{b5f});
\item intended use of the constants~$s_1,\dots,s_5$ is described by
  formulas~(\ref{d1}), (\ref{d2}), (\ref{d3}), (\ref{d4}), (\ref{d5}).
\end{itemize}

\section{Completions of programs A and B}\label{sec:comp}

Sentence~(\ref{a1f}) is not completable, because~$I$ and~$J$ are not
general variables~\cite[Section~3.1]{fan22}.  We rewrite it as
\beq
\ba r
\forall I J N X Y (\neg\neg q(I,J) \land 1\leq I \leq N \land 1 \leq J \leq N
\land N=n \land X=I \land Y=J \qquad\\
\to q(X,Y)).
\ea
\eeq{a1fcomp}
The set~$\Gamma_A$, consisting of sentences~(\ref{a2f}), (\ref{a3f}), (\ref{a4f}), (\ref{a5f}), and~(\ref{a1fcomp}), is completable, and its
completion is the conjunction of constraints~(\ref{a2f}), (\ref{a3f}),
(\ref{a4f}), (\ref{a5f}) and the completed definition of~$q$:
\beq
\ba l
\forall X Y (q(X,Y) \lrar\\ \hskip 1cm
\exists IJN
(\neg\neg q(I,J) \land 1\leq I \leq N \land 1 \leq J \leq N\land N=n
\land X=I \land Y=J)).
\ea
\eeq{acomp}
  Similarly,
the set~$\Gamma_B$, consisting of sentences~(\ref{b2f}), (\ref{b3f}),
(\ref{b4f}), (\ref{b5f}), and~(\ref{a1fcomp}), is completable, and its
completion is the conjunction of sentences~(\ref{b2f}), (\ref{b3f}),
(\ref{b4f}), (\ref{b5f}) and~(\ref{acomp}).

We will show how to derive the equivalence between the completion
of~$\Gamma_A$ and the completion of~$\Gamma_B$ in classical logic from
\begin{itemize}
\item \emph{KB}, a \emph{knowledge base} describing integers, precomputed
  terms, sets and cardinalities (see Section~\ref{sec:sets}),
\item \emph{UG}, the pair of \emph{user guide} formulas $\exists N(n=N)$
  and $n\geq 1$, expressing that~$n$ is a placehodler for a positive
  integer, and
\item \emph{SD}, the \emph{set definition} formulas~(\ref{d1}), (\ref{d2}),
  (\ref{d3}), (\ref{d4}), (\ref{d5}), which characterize the constants
  $s_1,\dots,s_5$.
\end{itemize}

\section{Plan of the proof}
 
As the first step, we will simplify the
completions of~$\Gamma_A$ and~$\Gamma_B$.  Formula~(\ref{acomp})
is equivalent to the implication
\beq
\ba l
\forall X Y (q(X,Y) \to \\\hskip 15mm
1 \leq X \leq n \land 1 \leq Y \leq n \land \exists I(X=I)
\land \exists J(Y=J) \land \exists N(N=n)).
\ea
 \eeq{iv}
Constraint (\ref{a2f}) can be rewritten as
\beq
\co(s_1)\geq n,
\eeq{a2fgeq}
\beq
\co(s_1)\leq n.
\eeq{a2fleq}
Constraint~(\ref{a3f}) can be rewritten as
$$\forall XY_1Y_2(q(X,Y_1) \land q(X,Y_2) \to Y_1 \geq Y_2),$$
and it implies
$$\forall XY_1Y_2(q(X,Y_1) \land q(X,Y_2) \to Y_2 \geq Y_1)$$
(substitute~$Y_2$,~$Y_1$ for~$Y_1$,~$Y_2$). It follows that~(\ref{a3f}) is
equivalent to
\beq
\forall XY_1Y_2(q(X,Y_1) \land q(X,Y_2) \to Y_1 = Y_2).
\eeq{a3fs}
Similarly,~(\ref{a4f}) is equivalent to
\beq
\forall X_1X_2Y(q(X_1,Y) \land q(X_2,Y) \to X_1 = X_2).
\eeq{a4fs}

Formula~(\ref{a5f})
is equivalent to
$$\forall I_1I_2J_1J_2
(q(I_1,J_1) \land q(I_2,J_2) \land |I_1-I_2|=|J_1-J_2| \to I_1 \geq I_2).$$
It implies
$$\forall I_1I_2J_1J_2
(q(I_1,J_1) \land q(I_2,J_2) \land |I_1-I_2|=|J_1-J_2| \to I_2 \geq I_1)
$$
(substitute~$I_2,\,I_1,\,J_2,\,J_1$ for $I_1,\,I_2,\,J_1,\,J_2$).
Hence it can be rewritten as
$$
\forall I_1I_2J_1J_2
(q(I_1,J_1) \land q(I_2,J_2) \land |I_1-I_2|=|J_1-J_2|\to I_1= I_2)
$$
and as
\beq
\forall I_1I_2J_1J_2
(q(I_1,J_1)\land q(I_2,J_2) \land |I_1-I_2|=|J_1-J_2|
\to (I_1,J_1) = (I_2,J_2)).
\eeq{a5fs}

Constraints~(\ref{b2f})~(\ref{b3f}),~(\ref{b4f}) and~(\ref{b5f})
can be rewritten as
\beq
\forall I(1\leq I\leq n \to \co(s_2(I)) \geq 1),
\eeq{b2fgeq}
\beq
\forall I(1\leq I\leq n \to \co(s_2(I))\leq 1),
\eeq{b2fleq}
\beq
\forall I(1\leq I\leq n \to \co(s_3(I))\geq 1),
\eeq{b3fgeq}
\beq
\forall I(1\leq I\leq n \to \co(s_3(I))\leq 1),
\eeq{b3fleq}
\beq
\forall KN(1\leq K\leq 2\times N-1  \land N=n \to \co(s_4(K,N))\leq 1),
\eeq{b4fs}
\beq
\forall KN(1\leq K\leq 2\times N-1 \land N=n \to \co(s_5(K))\leq 1).
\eeq{b5fs}

What we want to accomplish can be now described as follows:

\medskip\begin{center}
prove the equivalence\\
between the conjunction of formulas~(\ref{a2fgeq})--(\ref{a5fs})\\
and the conjunction of formulas~(\ref{b2fgeq})--(\ref{b5fs})\\
assuming \kbugsd\ and (\ref{iv}).
\end{center}\medskip

To this end, we will prove, under assumptions \kbugsd\ and (\ref{iv}),
a number of claims.

\medskip\emph{Claim 1:}
(\ref{a2fgeq}) follows from (\ref{b2fgeq}).

\medskip\emph{Claim 2:}
(\ref{a2fleq}) follows from~(\ref{b2fleq}).

\medskip\emph{Claim 3:}
(\ref{a3fs}) is equivalent to (\ref{b2fleq}).

\medskip\emph{Claim 4:}
(\ref{a4fs}) is equivalent to (\ref{b3fleq}).

\medskip\emph{Claim 5:}
(\ref{a5fs}) is equivalent to the conjunction of (\ref{b4fs}) and
(\ref{b5fs}).

\medskip\emph{Claim 6:}
(\ref{b2fgeq}) follows from  the conjunction of
(\ref{a2fgeq}) and (\ref{a3fs}).

\medskip\emph{Claim 7:}
(\ref{b3fgeq}) follows from  the conjunction of
(\ref{a2fgeq}) and (\ref{a4fs}).

\section{The knowledge base}\label{sec:sets}

The knowedge base \emph{KB} includes basic properties of integers and
arithmetic operations.  It includes also the formula
\beq
\forall X_1Y_1X_2Y_2((X_1,Y_1)=(X_2,Y_2) \to X_1=X_2 \land Y_1=Y_2),
\eeq{op}
the formula
\beq
\forall IJX(I\leq X\leq J \to \exists K(X=K)),
\eeq{contig}
which expresses that numerals are contiguous within precomputed terms,
and the universal closures of
formulas (\ref{exten})--(\ref{pigeon}) below, which express properties of
sets and cardinalities.

The axiom of extensionality
expresses that a set is determined by its extension:
\beq
\forall X(X\in S_1 \lrar X\in S_2) \to S_1=S_2.
\eeq{exten}

The axiom schema of comprehension expresses that
every property defines a set:
\beq
\exists S\forall X(X\in S \lrar F(X)),
\eeq{comprehension}
where~$F(X)$ is a formula without free occurrences of~$S$.

Cardinalities are nonnegative:
\beq
\co(S)=I \to I\geq 0.
\eeq{cardnn}

Comparing cardinalities with 1:
\begin{align}
&\fnt(S) \to (\co(S)\geq 1 \lrar \exists X(X\in S)),\label{geq1}\\
&\fnt(S) \to (\co(S)\leq 1 \lrar
\forall XY(X\in S \land Y\in S \to X=Y)).\label{leq1}
\end{align}

Cardinality of an interval:
\beq
I\leq J\land \forall X(X\in S \lrar I\leq X\leq J) \to \co(S) = J-I+1.
\eeq{cardint}

Cardinality of the union of disjoint sets:
\beq
\ba l
\co(S_1)=I \land \co(S_2)=J\\
\hskip 1cm
\land\; \neg\exists X(X\in S_1\land X\in S_2)\\
\hskip 1cm
\land\;\forall X(X\in S \lrar X\in S_1 \lor X\in S_2)\\
\hskip 5cm
\to \co(S)=I+J.
\ea\eeq{cardunion}

Cardinality of the Cartesian product:
\beq\ba l
\co(S_1)=I \land \co(S_2)=J\\
\hskip .5cm
\land\;\forall X(X\in S \lrar \exists Y_1Y_2(Y_1\in S_1 \land Y_2\in S_2
\land X=(Y_1,Y_2)))\\
\hskip 7.5cm
\to \co(S)=I\times J.
\ea\eeq{cardprod}

Finally, the pigeonhole principle:
\beq
\ba l
\forall X(X\in S_1 \to \exists Y (Y\in S_2 \land (X,Y)\in S))\\
\hskip .5cm
\land\;\forall X_1X_2Y((X_1,Y)\in S \land (X_2,Y)\in S \to X_1=X_2)
\land \fnt(S_2)\\
\hskip 6cm \to \fnt(S_1) \land \co(S_1)\leq \co(S_2).
\ea
\eeq{pigeon}

\begin{proposition}
The following properties of cardinalities are entailed by KB:
\beq
\ba l
\forall S_1S_2(S_1\subseteq S_2 \land \fnt(S_2)\\
\hskip 3cm
\to \fnt(S_1) \land \co(S_1)\leq \co(S_2)).
\ea
\eeq{subset}
\beq
\forall IJS(\forall X(X\in S \to I\leq X\leq J) \to \fnt(S)).
\eeq{intfinite}
\beq
\ba l
\forall S_1S_2S(\fnt(S_1) \land \fnt(S_2)
\land\forall X(X\in S \lrar X\in S_1 \lor X\in S_2)\\
\hskip 9.5cm
\to \fnt(S)).
\ea\eeq{unionfinite}
\beq
\forall S_1S_2(
S_1\subseteq S_2 \land \fnt(S_2) \land \co(S_1) \geq \co(S_2)\to S_1=S_2).
\eeq{propersubset}
\end{proposition}

\begin{proof}
% Formula~(\ref{intfinite}) is immediate from~(\ref{cardint}).
Formula~(\ref{subset}) is a special case of the pigeonhole
principle~(\ref{pigeon}): define~$S$ by the condition
$$\forall V(V\in S \lrar \exists X(V=(X,X))).$$
(Such~$S$ exists by comprehension~(\ref{comprehension}).)

To prove~(\ref{intfinite}), assume $\forall X(X\in S \to I\leq X\leq J)$
and consider the set~$S_1$ defined by the condition
$$\forall V(V\in S_1\lrar I\leq V\leq K),$$
where~$K$ is the greater of the numbers $I$, $J$.  Then $I\leq K$, and
by~(\ref{cardint}), $\fnt(S_1)$.  On the other hand, $S\subseteq S_1$.
Hence $\fnt(S)$ follows by~(\ref{subset}).

To prove~(\ref{unionfinite}), assume
  $$
  \fnt(S_1) \land \fnt(S_2)
  \land\forall V(V\in S \lrar V\in S_1 \lor V\in S_2),$$
  and consider the set~$S_3$ defined by the condition
  $$\forall V(V\in S_3 \lrar V\in S_1 \land V\not\in S_2).$$
Then
$$\neg\exists X(X\in S_2\land X\in S_3)$$
and
$$\forall X(X\in S \lrar X\in S_2 \lor X\in S_3).$$
On the other hand, from $\fnt(S_1)$ and~(\ref{subset}) we can derive
$\fnt(S_3)$,
Now $\fnt(S)$ follows by~(\ref{cardunion}) applied to~$S_2$,~$S_3$
as~$S_1$,~$S_2$.

% Formula~(\ref{prodfinite}) is immediate from~(\ref{cardprod}).

To prove~(\ref{propersubset}), assume
$$
S_1\subseteq S_2 \land \fnt(S_2) \land \co(S_1) \geq \co(S_2).
$$
From $S_1\subseteq S_2$ and $\fnt(S_2)$ we conclude by~(\ref{subset}) that
$\co(S_1) \leq \co(S_2)$; hence
%In view of the first conjunctive term, it is sufficient to show that
%$S_2 \subseteq S_1$;
%then \hbox{$S_1=S_2$} will follow by extensionality.
%From the last two conjunctive terms,
$$\co(S_1)=\co(S_2)=I$$
for some integer~$I$.  Consider the set~$S_3$ defined by the condition
$$\forall V(V\in S_3 \lrar V\in S_2 \land V\not\in S_1).$$
From $S_3\subseteq S_2$ we conclude by~(\ref{subset}) that
$\fnt(S_3)$, so that $\co(S_3)=J$ for some integer~$J$.  Hence,
by~(\ref{cardunion}), $I+J=I$,
so that~$J=0$. By~(\ref{leq1}), we conclude that $\neg\exists X(X\in S_3)$,
which implies $S_2\subseteq S_1$.  On the other hand, $S_1\subseteq S_2$;
by extensionality, it follows that $S_1=S_2$.
\end{proof}

\section{Proofs of Claims 1--7}

In this section, whenever a formula is asserted,
this is understood as the claim
that it is derivable in classical logic from assumptions \kbugsd\
and~(\ref{iv}).  Claims that one formula
implies another, or that two formulas are equivalent, are understood here in a
similar way.

\subsection{Two lemmas}

\begin{lemma}\label{lem:S0}
$\exists S_0(\forall V(V\in S_0\lrar 1\leq V\leq n)
  \land \fnt(S_0) \land \co(S_0)=n)$.
\end{lemma}
\begin{proof}
  The existence of~$S_0$ such that
  $\forall V(V\in S_0\lrar 1\leq V\leq n)$
  follows from comprehension~(\ref{comprehension}).  By \emph{UG}, $n=N\geq 1$
  for some integer~$N$.  By~(\ref{cardint}), it follows that
$\co(S_0) = N-1+1=N$, which implies both $\fnt(S_0)$ and $\co(S_0)=n)$.
\end{proof}
  
\begin{lemma}\label{lem:finite}
\strut
\begin{itemize}
\item[(i)]
  $\fnt(s_1)$.
\item[(ii)]
  $\forall I\, \fnt(s_2(I))$.
\item[(iii)]
  $\forall I\, \fnt(s_3(I))$.
\item[(iv)]
  $\forall KN\, \fnt(s_4(K,N))$.
\item[(v)]
$\forall N \fnt(s_5(N))$.
\end{itemize}
\end{lemma}
\begin{proof}
(i) Define~$S_1$ by the condition
$$
\forall V(V\in S_1 \lrar \exists XY(X\in S_0 \land Y\in S_0\land V=(X,Y))),
$$
where~$S_0$ is the set from Lemma~\ref{lem:S0},
By~(\ref{cardprod}), $\fnt(S_1)$.  On the other hand, by~(\ref{d1}),
$s_1\subseteq S_1$,
and $\fnt(s_1)$ follows by~(\ref{subset}).

(ii) Formulas~(\ref{d2}) and~(\ref{iv}) imply
$$\forall IY(Y\in s_2(I) \to 1\leq Y\leq n),$$
which is equivalent to
$\forall I(s_2(I)\subseteq S_0)$
for the set~$S_0$ from Lemma~\ref{lem:S0}.
Then $\forall I\,\fnt(s_2(I))$ follows by~(\ref{subset}).

For parts~(iii)--(v), the  reasoning is similar.
\end{proof}

\subsection{Proof of Claim~1}

By~(\ref{iv}), the last two conjunctive terms in the right-hand side
of~(\ref{d1}) can be dropped, so that
\beq
\forall V(V\in s_1 \lrar \exists IJ(V=(I,J)\land q(I,J))).
\eeq{d1s}

Assume~(\ref{b2fgeq}).  By~(\ref{geq1}) and Lemma~\ref{lem:finite}(ii),
$$
\forall I(1\leq I\leq n \to \exists X(X\in s_2(I))).
$$
Then, by~(\ref{d2}),
$$\forall I(1\leq I\leq n \to \exists X\, q(I,X)).$$
By~(\ref{iv}), it follows that
\beq
\forall I(1\leq I\leq n \to \exists J\, q(I,J)).
\eeq{p1a}

We want to apply pigeonhole principle~(\ref{pigeon}) to the set~$S_0$
from Lemma~\ref{lem:S0}, the set~$s_1$, and the set~$S$ defined
by the condition
$$\forall V(V\in S \lrar \exists XY(V=(X,(X,Y)))).$$
The three conjunctive terms in the antecedent of~(\ref{pigeon}) are verified
as follows.

To prove
\beq
\forall X(X\in S_0 \to \exists Y (Y\in s_1 \land (X,Y)\in S)),
\eeq{term1}
assume $X\in S_0$. Then $1\leq X\leq n$ and $X=I$ for some integer~$I$.
By~(\ref{p1a}), $q(I,J)$ for some~$J$.  By~(\ref{d1s}), it follows
that $(I,J)\in s_1$.  On
the other hand, $(I,(I,J))\in S$.  It follows that $(I,J)$ can be taken
as~$Y$ in the consequent of~(\ref{term1}).

To prove the conjunctive term
$$\forall X_1X_2Y((X_1,Y)\in S \land (X_2,Y)\in S \to X_1=X_2)$$
note that
$$\forall V_1V_2((V_1,V_2)\in S \lrar \exists Z(V_2=(V_1,Z))).$$
Assume $(X_1,Y)\in S \land (X_2,Y)\in S$.  Then there
exist~$Z_1$,~$Z_2$ such that $Y=(X_1,Z_1)$ and $Y=(X_2,Z_2)$.
Hence $X_1=X_2$.

The last conjunctive term $\fnt(s_1)$ is given by Lemma~\ref{lem:finite}.

By~(\ref{pigeon}), $\co(s_1)\geq \co(S_0)=n$.
\qed

\subsection{Proof of Claim~2}

Assume~(\ref{b2fleq}).  By~(\ref{leq1}) and Lemma~\ref{lem:finite}(ii),
$$\forall IXY(1\leq I\leq n \land X\in s_2(I) \land Y\in s_2(I) \to X=Y).$$
Hence, by~(\ref{d2}),
\beq
\forall IXY(1\leq I\leq n \land q(I,X) \land q(I,Y) \to X=Y).
\eeq{p2a}

We want to apply pigeonhole principle~(\ref{pigeon}) to the set~$s_1$, the
set~$S_0$ from Lemma~\ref{lem:S0}, and the set~$S$ defined by the condition
$$
\forall V(V\in S \lrar \exists IJ(V=((I,J),I)\land q(I,J))).
$$
The three conjunctive terms in the antecedent of~(\ref{pigeon}) are verified
as follows.

To prove
\beq
\forall X(X\in s_1 \to \exists Y (Y\in S_0 \land (X,Y)\in S)),
\eeq{term1a}
assume $X\in s_1$. Then, by~(\ref{d1}), there exist~$I$ and~$J$ such that
$$X=(I,J),\ q(I,J),\hbox{  and }1\leq I\leq n.$$
From~$q(I,J)$ we conclude that
$((I,J),I)\in S$.  Hence $I$ can be taken as~$Y$ in the consequent
of~(\ref{term1a}).

To prove the conjunctive term
$$\forall X_1X_2Y((X_1,Y)\in S \land (X_2,Y)\in S \to X_1=X_2)$$
note that
$$\forall V_1V_2((V_1,V_2)\in S \lrar
\exists IJ(V_1=(I,J) \land V_2=I \land q(I,J))).$$
Assume $(X_1,Y)\in S \land (X_2,Y)\in S$.  Then there
exist~$I_1$, $J_1$, $I_2$, $J_2$ such that
$$X_1=(I_1,J_1),\ X_2=(I_2,J_2),\ Y=I_1=I_2,\ q(I_1,J_1),\ q(I_2,J_2).$$
By~(\ref{iv}), $1\leq I_1=I_2\leq n$.
Hence, by~(\ref{p2a}), $J_1=J_2$, and consequently $X_1=X_2$.

The last conjunctive term $\fnt(S_0)$ holds by the choice of~$S_0$.

By~(\ref{pigeon}), $\co(s_1)\leq \co(S_0)=n$.
\qed

\subsection{Proofs of Claims~3 and~4}

The proof of Claim 3 uses the fact that, by~(\ref{leq1}) and
Lemma~\ref{lem:finite}(ii),
\beq
\forall I(\co(s_2(I))\leq 1 \lrar
\forall XY(X\in s_2(I) \land Y\in s_2(I)\to X=Y)).
\eeq{l}
By~(\ref{iv}), formula~(\ref{a3fs})
is equivalent to
$$
\forall I(1\leq I\leq n \to
\forall Y_1Y_2(q(I,Y_1) \land q(I,Y_2) \to Y_1 = Y_2)).
$$
By~(\ref{d2}), it can be further rewritten as
$$
\forall I(1\leq I\leq n \to
\forall Y_1Y_2(Y_1\in s_2(I) \land Y_2\in s_2(I) \to Y_1 = Y_2)).$$
By~(\ref{l}), this formula is equivalent to~(\ref{b2fleq}).
\qed

The proof of Claim~4 is similar.

\subsection{Proof of Claim~5}

\begin{lemma}\label{s4leq1}  For $i=4,5$,
  $$\ba l
  \forall KN(\co(s_i(K,N))\leq 1\\
\hskip 2cm
  \lrar \forall I_1I_2J_1J_2((I_1,J_1)\in s_i(K,N)
  \land (I_2,J_2)\in s_i(K,N)\\
  \hskip 8cm
  \to (I_1,J_1)=(I_2,J_2))).
  \ea$$
\end{lemma}

\begin{proof}
By~(\ref{leq1}) and Lemma~\ref{lem:finite}(iv),
\beq\ba l
\co(s_4(K,N))\leq 1\\
\hskip 2cm
\lrar \forall XY(X\in s_4(K,N) \land Y\in s_4(K,N) \to X=Y).
\ea\eeq{l2.1}
On the other hand, by~(\ref{d4}),
$$\forall X(X\in s_4(K,N) \to \exists IJ(X=(I,J))),$$
so that~(\ref{l2.1}) implies the assertion of the lemma for $i=4$.
For $i=5$ the proof is similar.
\end{proof}

\begin{lemma}\label{d4ij}
\strut\vskip 4mm
  
\begin{itemize}
\item[(i)]
$\forall IJKN((I,J)\in s_4(K,N) \lrar q(I,J) \land K=I-J+N)$,
\item[(ii)]
$\forall IJKN((I,J)\in s_5(K,N) \lrar q(I,J) \land K=I+J-1)$.
\end{itemize}
\end{lemma}
\begin{proof}
(i) By~(\ref{d4}), $(I,J)\in s_4(K,N)$ is equivalent to
$$\exists I_1J_1((I,J)=(I_1,J_1) \land q(I_1,J_1) \land K=I_1-J_1+N).$$
This formula is equivalent to the right-hand side of~(i) by~(\ref{op}).

The proof of~(ii) is similar, using~(\ref{d5}).
\end{proof}
\begin{lemma}\label{lem:at}
\strut\vskip 4mm
  
\begin{itemize}
\item[(i)]
Formula~(\ref{b4fs}) is equivalent to
\beq\forall I_1I_2J_1J_2
(q(I_1,J_1)\land q(I_2,J_2) \land I_1-J_1=I_2-J_2
\to (I_1,J_1) = (I_2,J_2)).
\eeq{a4up}
\item[(ii)]
Formula~(\ref{b5fs}) is equivalent to
\beq\forall I_1I_2J_1J_2
(q(I_1,J_1)\land q(I_2,J_2) \land I_1+J_1=I_2+J_2
\to (I_1,J_1) = (I_2,J_2)).
\eeq{a4down}
\end{itemize}
\end{lemma}
\begin{proof}
(i)~By Lemma~\ref{s4leq1},~(\ref{b4fs}) is equivalent to
$$
\ba l
\forall I_1I_2J_1J_2KN(1\leq K\leq 2\times N-1 \land N=n\\
\hskip 2.6cm
\land\; (I_1,J_1)\in s_4(K,N) \land (I_2,J_2)\in s_4(K,N)\\
\hskip 7cm
\to (I_1,J_1)=(I_2,J_2)).
\ea
$$
By Lemma~\ref{d4ij}, it can be further rewritten as
$$
\ba l
\forall I_1I_2J_1J_2KN(1\leq K\leq 2\times N-1 \land N=n\\
\hskip 2.6cm
\land\; q(I_1,J_1) \land K=I_1-J_1+N\\
\hskip 2.6cm
\land\; q(I_2,J_2) \land K=I_2-J_2+N\\
\hskip 5cm
\to (I_1,J_1)=(I_2,J_2)).
\ea
$$
By~(\ref{iv}), the conjunctive terms $N=n$ and $q(I_1,J_1)$ imply
$1\leq I_1 \leq N$ and  $1\leq J_1 \leq N$.  In combination
with the conjunctive term $K=I_1-J_1+N$, these inequalities imply
the conjunctive term $1\leq K\leq 2\times N-1$.  Hence that term can be
dropped, and the entire formula becomes
$$\ba l
\forall I_1I_2J_1J_2KN(N=n\land q(I_1,J_1) \land K=I_1-J_1+N\\
\hskip 3.7cm
\land\; q(I_2,J_2) \land K=I_2-J_2+N\\
\hskip 7cm
\to (I_1,J_1)=(I_2,J_2)).
\ea
$$
The conjunctive term $K=I_2-J_2+N$ here can be replaced by
$$I_1-J_1=I_2-J_2,$$
so that the formula can be further rewritten as
$$\ba l
\forall I_1I_2J_1J_2N(N=n\land q(I_1,J_1) \land \exists K(K=I_1-J_1+N)\\
\hskip 3.3cm
\land\; q(I_2,J_2) \land I_1-J_1=I_2-J_2\\
\hskip 7cm
\to (I_1,J_1)=(I_2,J_2)).
\ea
$$
After dropping the trivial conjunctive term $\exists K(K=I_1-J_1+N)$, the
formula becomes
$$\ba l
\forall I_1I_2J_1J_2(\exists N(N=n) \land q(I_1,J_1) \land
q(I_2,J_2) \land I_1-J_1=I_2-J_2\\
\hskip 8cm
\to (I_1,J_1)=(I_2,J_2)).
\ea
$$
By \emph{UG}, it is equivalent to~(\ref{a4up}).

The proof of part~(ii) is similar.
\end{proof}

\medskip\noindent\emph{Proof of Claim~5.}
Formula~(\ref{a5fs}) is clearly equivalent to the conjunction of~(\ref{a4up})
and~(\ref{a4down}), and, by Lemma~\ref{lem:at}, to the conjunction
of~(\ref{b4fs}) and~(\ref{b5fs}).
\qed

\subsection{Proofs of Claims~6 and~7}

To prove Claim~6, assume~(\ref{a2fgeq}) and (\ref{a3fs}).
We want to apply pigeonhole principle~(\ref{pigeon}) to~$s_1$
as~$S_1$ and to the sets~$S_2$ and~$S$ defined by the conditions
$$\ba l
\forall V(V\in S_2 \lrar \exists J\, q(V,J)),\\
\forall V(V\in S \lrar \exists IJ(V=((I,J),I)\land q(I,J))).
\ea$$
The three conjunctive terms in the antecedent of~(\ref{pigeon}) are verified
as follows.  To prove
$$\forall X(X\in s_1 \to \exists Y (Y\in S_2 \land (X,Y)\in S)),$$
assume $X\in s_1$.  By~(\ref{d1}), there exist~$I$,~$J$ such that
$X=(I,J)$ and $q(I,J)$; we can take~$I$ as~$Y$.

To prove
$$\forall X_1X_2Y((X_1,Y)\in S \land (X_2,Y)\in S \to X_1=X_2)$$
assume $(X_1,Y)\in S \land (X_2,Y)\in S$.  As in the proof of Claim~2,
there exist~$I_1$, $J_1$, $I_2$, $J_2$ such that
$$X_1=(I_1,J_1),\ X_2=(I_2,J_2),\ Y=I_1=I_2,\ q(I_1,J_1),\ q(I_2,J_2).$$
By~(\ref{a3fs}), it follows that $J_1=J_2$.  Hence $X_1=X_2$.

To prove the third conjunctive term $\fnt(S_2)$, note that $S_2\subseteq S_0$
for the set~$S_0$ from Lemma~\ref{lem:S0}; $\fnt(S_2)$ follows
by~(\ref{subset}).

By~(\ref{pigeon}), $\co(s_1)\leq \co(S_2)$.  By~(\ref{a2fgeq}), it
follows that $\co(S_2)\geq n=\co(S_0)$.  Hence, by~(\ref{propersubset}),
$S_2=S_0$, so that, for every $I$ such that $1\leq I\leq n$, $I\in S_2$.
This condition can be rewritten as $\exists J\,q(I,J)$.  By~(\ref{d2}),
this formula is equivalent to $\exists J\,(J\in s_2(I))$.  By~(\ref{geq1}),
it can be further rewritten as $\co(s_2(I))\geq 1$.

\section{Conclusion}

We proved that the equivalence between the completions of the
programs $A$ and $B$ can be derived from the assumptions \kbugsd.
Future work will include establishing a general theorem showing that
the existence of such a derivation implies the external equivalence of
the programs.  That theorem will be somewhat similar to the theorem on
completion and tightness of
programs with non-recursive aggregates \cite[Section~7.1]{fan24b}.

It would be interesting also to implement a version of {\sc anthem}
\cite{fan25} capable of verifying the derivation outlined in this note.

\bibliographystyle{named}
\bibliography{bib}
\end{adjustwidth}
\end{document}

